% Part: first-order-logic
% Chapter: syntax-and-semantics
% Section: covered-structures

\documentclass[../../../include/open-logic-section]{subfiles}

\begin{document}

\olfileid{fol}{syn}{cov}

\olsection{Gedekte \printtoken{P}{structure} voor talen van de eerste orde
(eerste-ordetalen)}

\begin{explain}
Een term heet \emph{gesloten} als er geen !!{variable}s in staan.
\end{explain}

\begin{defn}[De !!^{value} van een gesloten term]
Neem een gesloten term $t$ van de taal~$\Lang L$ en een
!!{structure}~$\Struct M$ voor~$\Lang L$. De
\emph{!!{value}}~$\Value{t}{M}$ is het element van het domein dat de
term binnen deze structuur aanwijst. Je bepaalt die waarde zo:
\begin{enumerate}
\item Bestaat $t$ alleen uit de !!{constant}~$c$, dan is
  $\Value{c}{M} = \Assign{c}{M}$.
\item Heeft $t$ de vorm $\Atom{f}{t_1, \ldots, t_n}$, bepaal dan eerst
  de waarden van de termen binnen de haakjes. Pas daarna de functie toe:
  \[
  \Value{t}{M} = \Assign{f}{M}(\Value{t_1}{M}, \ldots,
  \Value{t_n}{M}).
  \]
\end{enumerate}
\end{defn}

\begin{defn}[Gedekte !!{structure}]
Een !!{structure} heet \emph{gedekt} als elk element van het domein de
!!{value} van een gesloten term is. Je kunt dan dus elk domeinelement
met minstens één term zonder variabelen aanwijzen.
\end{defn}

\begin{ex}
Neem de taal~$\Lang L$ met de !!{constant}s $\Obj{zero}$,
$\Obj{one}$, $\Obj{two}$, \dots. Deze taal bevat ook het binaire
!!{predicate}~$<$ en de binaire !!{function}s $+$ en $\times$.
Neem als !!a{structure}~$\Struct M$ voor~$\Lang L$ de volgende. Het
domein is $\Domain M = \{0, 1, 2, \ldots \}$. De structuur legt het zo
vast: $\Assign{\Obj{zero}}{M} = 0$, $\Assign{\Obj{one}}{M} = 1$,
$\Assign{\Obj{two}}{M} = 2$, enzovoort. Voor het binaire
relatiesymbool $<$ bevat de verzameling $\Assign{<}{M}$ precies de
paren $\tuple{c_1, c_2} \in \Domain{M}^2$ waarvoor $c_1$ kleiner is
dan~$c_2$. Zo geldt $\tuple{1, 3} \in \Assign{<}{M}$; daarentegen
geldt $\tuple{2, 2} \notin \Assign{<}{M}$. Bij de
binaire !!{function} $+$ is $\Assign{+}{M}$ de gewone optelling. Zo
levert $\Assign{+}{M}(2,3)$ als uitkomst~$5$. Voor de binaire
!!{function}~$\times$ gebruiken we op dezelfde manier de gewone
vermenigvuldiging. Daarom is de !!{value} van $\Obj{four}$ gewoon~$4$.
De !!{value} van $\times(\Obj{two},
+(\Obj{three},\Obj{zero}))$ (of, met het bewerkingsteken tussen de
getallen: $\Obj{two} \times (\Obj{three} + \Obj{zero})$) is
\begin{multline*}
\Value{\times(\Obj{two}, +(\Obj{three},\Obj{zero}))}{M} =\\
\begin{aligned}
& =\Assign{\times}{M}(\Value{\Obj{two}}{M}, \Value{+(\Obj{three}, \Obj{zero})}{M})\\
& = \Assign{\times}{M}(\Value{\Obj{two}}{M}, \Assign{+}{M}(\Value{\Obj{three}}{M},
\Value{\Obj{zero}}{M})) \\
& = \Assign{\times}{M}(\Assign{\Obj{two}}{M}, \Assign{+}{M}(\Assign{\Obj{three}}{M},
\Assign{\Obj{zero}}{M})) \\
& = \Assign{\times}{M}(2, \Assign{+}{M}(3, 0)) \\
& = \Assign{\times}{M}(2, 3) \\
& = 6
\end{aligned}
\end{multline*}
\end{ex}

\begin{prob}
Is $\Struct N$, het standaardmodel van de rekenkunde, gedekt? Leg uit
waarom wel of niet.
\end{prob}

\end{document}
